\documentstyle {article}



\title{
A flexible tool to experiment with distributed symbolic computation
}

\author{
        St\'ephane Dalmas and Marc Ga\"etano
        \thanks{also I3S, CNRS URA 1376,Universit\'e 
                          de Nice Sophia--Antipolis}\\
        {\small INRIA Sophia Antipolis, projet SAFIR} \\
        {\small 2004 route des lucioles }
         {\small BP 93} \\
         {\small 06902 Sophia-Antipolis {\sc CEDEX}} \\
         {\small France } \\
         {\small 
          \begin{tabular} {ll}
            E-mail:& stephane.dalmas@sophia.inria.fr\\
                   & marc.gaetano@sophia.inria.fr
          \end{tabular} 
         }
}



\begin{document} 

\maketitle

\begin{abstract}
In this paper we describe the {\em Central Control\/} and how it can be
applied to easily make various experiments in distributed computation. 
The {\em Central Control\/} is a software component that has been designed to
be the kernel of  environments for scientific computations. 
Such an environment can offer a common and concurrent access to many tools
needed by the scientist and the engineer: general purpose and specialized
computer algebra systems, visualization tools, links with numerical libraries
and tools to manipulate numerical programs {\sl etc}.

The current version of the {\em Central Control\/} is a {\em Scheme\/}
interpreter extended with new primitive types and procedures. 
To insure flexibility and to ease experimentations, there are only a few new
essential procedures to provide the basic facilities (starting or connecting
to servers, sending and receiving computations) together with a {\em Scheme}
library that implements higher-level operations.

We will first describe the main features of the {\em Central Control\/}. We
will then give examples of how it can be applied to problems that need or can
benefit from distribution by showing examples of simple {\em Scheme}
functions.  

\end{abstract} 



\end{document} 
